\section{A repository reading that follows the same method}
We now return to the repository that motivated the course. Its documentation distinguishes manuscripts and verification scopes; a repository listing is not, by itself, a guarantee about every claim in every paper~\cite{oairepo}. The appropriate unit of reading is an exact theorem together with its definitions and dependencies, not a title or an impressive final line.

For Family~173, the statement concerns a nonempty finite oriented graph and a vertex whose second out-neighborhood has at least as many vertices as its first. Chapter~7 supplies the mathematical vocabulary; Chapter~16 explains the exclusions in second-neighborhood membership~\cite{oai173}. In the formal signature, \code{Fintype V} supplies finiteness, \code{Nonempty V} rules out the empty graph, and \code{DecidableEq V} provides a way to decide equality for the finite constructions. The edge relation has type \code{V -> V -> Prop}.

The entry point imports modules for reduction, pruning, and an extremal contradiction. Its proof assumes a counterexample, obtains a transformed finite oriented relation from a reduction theorem, and applies the pruning and extremal results to obtain a contradiction~\cite{oaimain}. The reduction's output is a relation on pairs of finite indices with additional properties, including positive indegree and a strict subset-growth condition~\cite{oaireduction}. Understanding those property names requires opening their definitions; translating the names into plausible English is not enough.

A responsible reading notebook therefore has three different entries: the final statement is decoded; the logical role of the reduction is identified; the mathematical proof of the reduction and the meaning of its additional predicates remain tasks until their source has been read. The full repository argument is not reproduced or independently verified in this book. We have prepared the vocabulary and method for entering it, not substituted a short dependency sketch for that argument.

The comparison configuration identifies the challenge module, the solution module, the theorem name, and permitted axioms~\cite{oaiconfig}. A placeholder in a challenge template is not equivalent to a placeholder in the submitted solution's proof. Conversely, a theorem name in a solution file is not sufficient evidence that it matches the challenge. The configuration and the actual checking procedure matter.

For Family~179, Chapter~14 lets you decode the matrix and correlation conditions and prove some necessary restrictions yourself. The repository's scope note reports a stronger classification and distinguishes its even-length Barker statement from odd-length claims outside that selected formalization~\cite{oai179}. You can now ask a precise next question: which lemma rules out the candidate orders $4u^2$ with $u>1$ that our elementary argument still permits? That is a mathematically informed question, unlike asking an assistant simply to ``explain the whole repository.''

\begin{lab}{Request background, then test whether it is sufficient}
Begin with Theorem~\ref{thm:subsequence} while temporarily hiding its proof. Ask the assistant for a dependency map using only the knowledge from Chapter~12. Reject a response that merely lists ``compactness, topology, sequences.'' Require an example, a nonexample, and the exact role of every proposed concept. After reconstructing the proof, close the explanation and answer: why must the sets of indices be nested, why must indices increase, and why is boundedness not enough? A map that leaves one of these questions unanswered is not yet sufficient.

Then open one of the two repository scope notes and its named entry point. Produce a one-page reading record distinguishing what you can now reconstruct from what is still borrowed. Use the source file names and actual theorem identifiers. Do not report a repository build or formal verification unless you have performed it and retained the result.
\end{lab}
